\chapter{はじめに}
\label{ch:introduction}

\section{本稿の位置付け}

本稿は，千葉工業大学 未来ロボティクス学科 2026 年度「ハッカソン 1」の\textbf{事前課題}として提出するものである。事前課題の要件は「チームで協議するための技術的・社会的背景，既存のシステム，自身のアイディアを参考文献とともに検討すること」であり，本稿はこれに対応する。

\section{チームテーマと筆者の担当}

チーム 23 のテーマは \textbf{Arduino UNO R4 WiFi 5 台による「Arduino オーケストラ」} である。指揮者役のマイコン 1 台と楽器役のマイコン 4 台が Wi-Fi（UDP）で連携し，音符データを PC（Processing）へ送って金管楽器の音色を合成・再生する構成とする。

筆者（塩澤）は \textbf{指揮者マイコン担当}であり，本稿では指揮者マイコンの入力となる「IMU（慣性計測装置）を用いた指揮ジェスチャ認識」を中心に議論を進める。

\section{本稿の構成}

本稿は以下の構成を取る。

\begin{itemize}
  \item 第\ref{ch:background}章 \textbf{背景}: 本テーマに関連する社会的背景・技術的背景を整理する。
  \item 第\ref{ch:existing_systems}章 \textbf{既存システム}: ジェスチャ指揮・分散演奏・合奏民主化の 3 軸で既存事例を比較し，本提案の差別化点を示す。
  \item 第\ref{ch:idea}章 \textbf{自身のアイディア}: Arduino オーケストラの具体構成，特に指揮者マイコン部の処理方針と拡張性を述べる。
  \item 第\ref{ch:conclusion}章 \textbf{まとめと今後の検討事項}: チームで協議したい事項と次週までの個人タスクを整理する。
\end{itemize}
